\subsection{Rancangan Format Data Lokal}\label{rancangan-solusi-format-data-lokal}

Untuk menangani mekanisme \textit{retry mechanism}, klien pasien menyimpan
setiap \textit{batch} PGHD dalam \textit{local database} dengan struktur pada
Tabel~\ref{tab:field-batch-pghd-local}.

\begin{table}[htbp]
    \centering
    \caption{Struktur Data Lokal Batch PGHD}\label{tab:field-batch-pghd-local}
    \begin{styledtable}{>{\raggedright\arraybackslash}p{0.3\textwidth} >{\raggedright\arraybackslash}p{0.18\textwidth} >{\raggedright\arraybackslash}X}
        \textbf{Field} & \textbf{Tipe} & \textbf{Deskripsi} \\
        \texttt{batch\_id} & UUID & Identitas unik batch, sama dengan field pada skema JSON (Tabel~\ref{tab:field-batch-json-pghd}). \\
        \texttt{enc\_pghd} & String (Base64) & Ciphertext hasil enkripsi (AES-GCM) dari payload PGHD\@. \\
        \texttt{enc\_aes\_key\_nonce} & String (Base64) & Kunci AES dan nonce yang dienkripsi untuk kebutuhan PRE nantinya\@. \\
        \texttt{capsule} & String (Base64) & Artefak kriptografis untuk kebutuhan PRE nantinya. \\
        \texttt{h\_cipher} & String (Hex) & Hash SHA-256 dari \texttt{enc\_pghd}. \\
        \texttt{pghd\_outer\_signature} & String (Base64) & Tanda tangan digital dari \texttt{h\_cipher}. \\
        \texttt{status} & Enum & Status pengiriman: \texttt{pending} (belum dikirim), \texttt{sent} (berhasil), \texttt{failed} (gagal sementara), \texttt{permanent\_failure} (gagal permanen). \\
        \texttt{retry\_count} & Integer & Jumlah percobaan pengiriman (default 0, maksimal 3). \\
        \texttt{created\_at} & Timestamp & Waktu batch dibuat (Unix timestamp). \\
        \texttt{last\_attempt} & Timestamp & Waktu percobaan pengiriman terakhir (Unix timestamp). \\
    \end{styledtable}
\end{table}

Sebelum PGHD dibentuk menjadi batch dan dienkripsi untuk dikirimkan ke sistem,
aplikasi pasien menyimpan data PGHD dalam bentuk lokal raw seperti dapat
dilihat pada Tabel~\ref{tab:field-pghd-raw-local}. Data ini digunakan agar
pasien dapat melihat kembali data kesehatannya sendiri, termasuk sumber data
dan status sinkronisasinya.

\begin{table}[H]
    \centering
    \caption{Format Data PGHD Lokal Raw}\label{tab:field-pghd-raw-local}
    \begin{styledtable}{>{\raggedright\arraybackslash}p{0.3\textwidth} >{\raggedright\arraybackslash}p{0.18\textwidth} >{\raggedright\arraybackslash}X}
        \textbf{Field} & \textbf{Tipe} & \textbf{Keterangan} \\
        \texttt{record\_id} & String & Identitas unik data PGHD lokal. \\
        \texttt{patient\_id} & String & Identitas pasien pada aplikasi. \\
        \texttt{type} & String & Jenis data kesehatan, misalnya langkah, detak jantung, atau kalori. \\
        \texttt{value} & Number & Nilai pengukuran. \\
        \texttt{unit} & String & Satuan pengukuran. \\
        \texttt{timestamp} & DateTime & Waktu pengukuran. \\
        \texttt{source} & String & Sumber data, misalnya Health Connect atau sensor perangkat. \\
        \texttt{source\_package} & String & Nama paket aplikasi sumber data apabila berasal dari Health Connect. \\
        \texttt{sync\_status} & Enum & Status sinkronisasi. \\
        \texttt{created\_at} & DateTime & Waktu data disimpan di perangkat. \\
    \end{styledtable}
\end{table}